\documentclass[11pt,a4paper]{article}

\usepackage[margin=25mm]{geometry}
\usepackage[T1]{fontenc}
\usepackage[utf8]{inputenc}
\usepackage{lmodern}
\usepackage{enumitem}
\usepackage{booktabs}
\usepackage{amsmath}
\usepackage[table]{xcolor}
\usepackage{colortbl} % \arrayrulecolor
\usepackage{float} % for [H] to keep tables from floating in the response letter
\usepackage{hyperref}

\setlength{\parskip}{0.6em}
\setlength{\parindent}{0pt}

\hypersetup{
  colorlinks=true,
  linkcolor=black,
  urlcolor=blue!60!black,
  citecolor=black,
}

\definecolor{responsegreen}{HTML}{1B6E2B}

\newenvironment{commentblock}{%
  \par\smallskip\noindent\textbf{Comment:}\par
  \begin{quote}\itshape\color{black}
}{%
  \end{quote}\smallskip
}

\newenvironment{responseblock}{%
  \par\smallskip\noindent\textbf{Response:}\par
  \begingroup\color{responsegreen}
}{%
  \endgroup\smallskip
}

\title{Response to Reviewers}
\author{%
  \parbox{\textwidth}{%
    \textbf{Manuscript:} R4Tun: LLM-guided adaptive segmental tunnel lining segmentation in point clouds\\[0.4em]
    \textbf{Manuscript number:} AUTCON-D-26-02407
  }%
}
\date{}

\begin{document}
\maketitle

We are grateful to the Associate Editor and reviewers for their time in reviewing the manuscript and for their useful and insightful feedback. We have responded to the queries individually below and modified the manuscript as appropriate.

% \medskip
% \noindent\textbf{Summary of major revisions.} The revised manuscript:
% \begin{itemize}[nosep]
%   \item reframes R4Tun as a \emph{mechanism-contribution} study (bounded, cross-LLM, zero-label parameter adaptation) rather than a deployment-ready segmentation method;
%   \item adds an error-analysis subsection (Section~4.6; Figs.~11--12; Tables~9--10) decomposing false negatives, false positives, and class swaps, and quantifying four structural limits of the fixed SAM4Tun labelling substrate;
%   \item adds external positioning against supervised and expert-tuned tunnel-segmentation methods (Section~5.2; Table~11);
%   \item reports a runtime trade-off against the base SAM4Tun processing time (Section~4.4; Table~7);
%   \item adds a repeatability appendix (Appendix~I; Table~25);
%   \item expands limitations in the Conclusion, distinct from future-work directions.
% \end{itemize}
% \noindent All changes are highlighted in the revised manuscript PDF.

\section*{Associate Editor}

\begin{commentblock}
Avoid orphan headings without context: these are headers that are immediately followed by the next header, without any text in between (line 85--86, 169--170, \ldots). At least add a brief introduction to explain what will be covered in the section.
\end{commentblock}
\begin{responseblock}
Short introductory paragraphs have been added at the opening of Sections~2--5 (Related work, Methodology, Results, and Discussion), each summarising what the section covers.
\end{responseblock}

\begin{commentblock}
Fix reference on line 191. Also in other places, the reference to the Appendix is missing. All references in the PDF need to resolve to a number. The Appendix can be directly included in the manuscript after the references.
\end{commentblock}
\begin{responseblock}
Unresolved appendix references have been corrected. All appendix cross-references now resolve to numbered appendices (A--I). The Appendix is included directly in the manuscript after the References.
\end{responseblock}

\begin{commentblock}
Limitations must be included and discussed in the conclusion. Limitations are also not necessarily the same thing as Future Work, as some limitations are simply inherently part of the proposed approach (and thus need to be highlighted in a conclusion as well).
\end{commentblock}
\begin{responseblock}
Section~6 (Conclusions) now contains a dedicated limitations paragraph, separate from the future-work directions that follow. The limitations discussed include:
\begin{itemize}[nosep]
  \item evaluation on SAM4Tun and Seg2Tunnel only;
  \item single expert-tuned reference anchor;
  \item inherited SAM4Tun operator assumptions and failure modes;
  \item dependence on offline bounds and knowledge quality;
  \item single-pass adaptation without iterative self-correction.
\end{itemize}
\end{responseblock}

\begin{commentblock}
Do a thorough check of all references such that they follow the author guidelines and include full bibliographic details. Author names, journal or book titles, chapter or article titles, year of publication, volume numbers, article numbers or pagination must be included, where applicable. Several names are mistyped, likely due to improper handling in \LaTeX. Include DOIs for articles and ISBNs for book volumes. Check whether arXiv preprint articles have not been published by now, and try to replace them in that case, such that you can rely more on peer-reviewed and validated sources rather than pre-prints. Be careful with capitalization in titles as well, when using \LaTeX; many acronyms are now in lower-case. When citing an article with 6+ authors, only the first six authors need to be listed, and the remaining ones can be replaced by et al. More guidance on this can be found in the author guide of the journal, and the examples it refers to.
\end{commentblock}
\begin{responseblock}
The reference list has been checked for author names, article titles, journal titles, volume/article numbers, page ranges, DOIs, capitalisation, and published versions of preprints where available. \LaTeX{} capitalisation errors for acronyms (e.g., SAM, LLM, GPR) have been corrected, long author lists truncated to six authors plus ``et al.'' where required, and DOIs added for journal articles.
\end{responseblock}

\section*{Reviewer 1}

\begin{commentblock}
The manuscript is generally well revised and significantly improved. One minor suggestion is to further strengthen the discussion of recent studies on intelligent tunnel inspection, adaptive parameter tuning, and AI-assisted point-cloud analysis for underground infrastructure. In particular, the authors may consider citing several recent related works to better position the contribution of R4Tun within the latest development of LLM-/AI-enabled tunnel analysis and infrastructure inspection research:
\begin{itemize}[nosep]
  \item DOI: \url{10.1016/j.aei.2026.104774}
  \item DOI: \url{10.1016/j.autcon.2026.106769}
  \item DOI: \url{10.1016/j.tust.2024.105998}
  \item DOI: \url{10.1016/j.measurement.2023.113903}
\end{itemize}
\end{commentblock}
\begin{responseblock}
We thank the reviewer for this positive assessment. We have added six recent references to Related Work and Discussion:
\begin{itemize}[nosep]
  \item Yue et al.\ (TUST~2024; DOI \url{https://doi.org/10.1016/j.tust.2024.105998}) on shield-tunnel lining damage mechanisms;
  \item Yue et al.\ (Measurement~2024; DOI \url{https://doi.org/10.1016/j.measurement.2023.113903}) on GPR-based tunnel-lining defect recognition;
  \item Wang et al.\ (AUTCON~2026; DOI \url{https://doi.org/10.1016/j.autcon.2026.106908}) on point-cloud segmentation for geometric twins of shield metro tunnels;
  \item Wang et al.\ (Applied Sciences~2025; DOI \url{https://doi.org/10.3390/app15062926}) on automated tunnel point cloud segmentation and extraction;
  \item Yang et al.\ (AUTCON~2024; DOI \url{https://doi.org/10.1016/j.autcon.2024.105818}) on YOLO-based crack segmentation and measurement in metro tunnels;
  \item Liang et al.\ (TUST~2026; DOI \url{https://doi.org/10.1016/j.tust.2026.107705}) on multi-defect inspection of tunnel linings with YOLO+SAHI.
\end{itemize}
These citations better position R4Tun within recent AI-enabled tunnel and infrastructure inspection research.
\end{responseblock}

\begin{commentblock}
The proposed framework currently evaluates adaptation quality mainly through segmentation metrics (mIoU and OA). It would be helpful if the authors could briefly discuss whether the adapted parameters remain stable across repeated runs or different API calls, especially considering that LLM-based reasoning may still introduce slight stochasticity even under low-temperature settings. A short discussion on reproducibility would further strengthen the practical applicability of the framework.
\end{commentblock}
\begin{responseblock}
A repeatability check under the full $m+s+k$ setting (temperature~0) indicates limited stochastic drift. Across 30 tunnels, each LLM was re-inferred twice and compared between run~1 and run~2: on average, 90.9\% of the 18 critical parameters remain unchanged, and the mean absolute inter-run $|\Delta\text{mIoU}|$ is 0.029 (median 0.000), as reported in Table~20 (Appendix~7).
\end{responseblock}

\begin{commentblock}
The paper would benefit from a short clarification regarding the computational overhead introduced by the multi-agent adaptation process. Although runtime and API costs are reported, a brief comparison between the adaptation time and the original SAM4Tun inference time would help readers better understand the trade-off between adaptability and efficiency in practical deployment scenarios.
\end{commentblock}
\begin{responseblock}
We have added a direct runtime comparison in Section~4.4 (Runtime trade-off; Table~7). We separated the base SAM4Tun processing time (235~s per tunnel) from the additional LLM adaptation time under $m+s+k$:

{%
\begin{table}[H]
  \centering
  \small
  \color{responsegreen}\arrayrulecolor{responsegreen}
  \label{tab:runtime-tradeoff}
  \begin{tabular}{@{}lccc@{}}
    \toprule
    LLM & Extra time & Total time & Mean $\Delta\text{mIoU}$ \\
    \midrule
    Gemini-3-Flash & +96~s (0.4$\times$)  & $\sim$331~s & +0.166 \\
    Opus-4.6       & +140~s (0.6$\times$) & $\sim$375~s & +0.192 \\
    GPT-5.4        & +307~s (1.3$\times$) & $\sim$542~s & +0.171 \\
    \bottomrule
  \end{tabular}
\end{table}
}%
\end{responseblock}

\begin{commentblock}
In Section~5.2, the limitations are discussed clearly. As a minor improvement, the authors could also briefly comment on how the proposed framework might generalize to other point-cloud processing pipelines beyond SAM4Tun, for example whether the context-design strategy (memory-state-knowledge) is expected to remain transferable even when the downstream segmentation backbone changes.
\end{commentblock}
\begin{responseblock}
Section~5.3 (Limitations) acknowledges that transferability to other pipelines is conceptualised without empirical validation. The $m+s+k$ design can apply to other parameter-sensitive multi-stage point-cloud pipelines if four conditions hold:
\begin{itemize}[nosep]
  \item bounded tunable stage parameters;
  \item reference configurations with comparable target characterisation;
  \item summarisable intermediate state;
  \item encodable stage-specific knowledge.
\end{itemize}
\end{responseblock}

\section*{Reviewer 2}

\begin{commentblock}
I would like to commend the authors for their efforts in revising this manuscript. The authors have clearly taken the previous reviewers' comments. The expansion of the dataset to 30 tunnel subsets, the introduction of a deterministic non-LLM rule-based baseline, the rigorous cumulative ablation studies, and the inclusion of comprehensive statistical analyses have significantly elevated the methodological rigor of the paper. The concept of utilizing LLMs as white-box reasoning agents to adapt expert-designed pipelines via structured context is interesting.
\end{commentblock}
\begin{responseblock}
We thank the reviewer for this positive assessment and for highlighting the concept of our approach.
\end{responseblock}

\begin{commentblock}
However, despite the exhaustive experimental design, I have a fundamental and critical concern regarding the performance of the proposed system. The empirical results, as currently presented, fall drastically short of the standards required for practical deployment in construction automation, which severely undermines the core premise of the paper.
\end{commentblock}
\begin{responseblock}
We agree that the absolute metrics are below deployment thresholds for autonomous inspection. We have therefore reframed the study as an exploration of a controlled, cross-LLM, zero-label, no-retraining, self-adaptation \emph{mechanism}, rather than a performance-competitive segmentation method. Operating in a single forward pass ($\sim$331~s to $\sim$542~s per tunnel), R4Tun introduces explicit LLM reasoning into an expert-designed pipeline and serves as a methodological foundation for future development.
\end{responseblock}

\begin{commentblock}
\textbf{Massive Misclassification:} An OA of $\sim$53\% implies that nearly half of the point cloud is incorrectly classified. An mIoU of $\sim$33\% indicates a drastic lack of overlap between the predicted segment masks and the ground truth.
\end{commentblock}
\begin{responseblock}
We agree that these absolute figures are low. Section~4.6 (Error analysis) therefore decomposes the residual error into (i) errors that are controllable via LLM-selected parameters and (ii) errors that are structural to the fixed SAM4Tun labelling rule.

(i) On regular tunnels, tuning the keystone placement and related parameters largely eliminates under-segmentation, reducing false negatives from 21\% to 2\% of ground-truth points. On complex tunnels, tuning radius- and geometry-related parameters recovers under-segmented blocks (false negatives 66\% $\rightarrow$ 17\%). 

(ii) However, most remaining mistakes are class swaps: because SAM4Tun assigns labels via a fixed segment-offset template and a hard-coded segment order in the open-sourced implementation, swapped labels persist even when the correct points are recovered (26\% $\rightarrow$ 21\% on regular tunnels; rising to 34\% on complex tunnels). 

False positives remain small ($\leq$6\%) throughout. In short, adaptation corrects the dominant ``missing segment'' failure mode, while the remaining error is dominated by the fixed positional template rather than by LLM parameter selection.
\end{responseblock}

\begin{commentblock}
\textbf{Failure in Downstream Tasks:} In the context of automated tunnel inspection (e.g., structural health monitoring, deformation analysis, water leakage detection, or Scan-to-BIM workflows), a segmentation pipeline with a $\sim$50\% error rate is practically unusable. The downstream geometric and structural analyses would be entirely compromised by the massive amount of noise, missing segments, and misidentified structural components.
\end{commentblock}
\begin{responseblock}
We agree. The revised manuscript states explicitly in the Introduction, Discussion, and Conclusion that the current absolute accuracy does not yet sufficiently support autonomous final inspection or downstream tasks such as structural health monitoring, deformation analysis, water-leakage detection, or Scan-to-BIM without human verification.
\end{responseblock}

\begin{commentblock}
\textbf{The Low Baseline Fallacy:} The authors celebrate a $>$100\% relative improvement, but this is a classic statistical illusion caused by a catastrophically failing baseline. The static SAM4Tun baseline completely collapses on off-reference tunnels (overall mIoU = 0.15; complex tunnels = 0.04). Recovering from a state of total failure to a state of mediocre/poor performance (mIoU 0.34) does not validate the framework for real-world engineering applications.
\end{commentblock}
\begin{responseblock}
We agree that headline relative-gain framing can be misleading when the baseline is near zero. We have removed celebratory relative-improvement language (e.g., ``$>$100\% relative improvement'') from the Abstract, Discussion, and Conclusion. Discussion now leads with absolute mIoU and OA; paired absolute deltas (e.g., 0.15 $\rightarrow$ 0.32--0.34) are retained only where they aid interpretation of the adaptation effect. We do not present recovery from baseline collapse as evidence of deployment readiness.
\end{responseblock}

\begin{commentblock}
\textbf{Poor Performance on Regular Tunnels:} Most concerningly, even on the regular tunnel subsets---which closely resemble the expert-tuned reference configuration---the mIoU only reaches $\sim$0.50--0.54. If an LLM, armed with explicit memory, intermediate state, and expert knowledge documents, cannot achieve a high degree of accuracy (e.g., mIoU $>$ 0.80) on in-distribution or near-reference data, it raises serious doubts about either the LLM's actual reasoning capability in this domain or the fundamental brittleness of the underlying SAM4Tun pipeline.
\end{commentblock}
\begin{responseblock}
We agree this is a fair concern. Section~4.6 shows that on regular tunnels the LLM recovers detection (false negatives drop from 21\% to 2\%) but cannot eliminate class swaps (26\% $\rightarrow$ 21\%), which trace to the fixed segment-offset and segment-order template rather than incorrect parameter choices. The near-reference mIoU of 0.50--0.54 therefore reflects a \emph{substrate ceiling} of the single-reference SAM4Tun labelling rule, not a failure of LLM reasoning over structured context. This is consistent with the Non-LLM vs.\ LLM comparison: on regular tunnels, deterministic rules do not beat the static baseline, whereas the LLM ($m+s+k$) does.
\end{responseblock}

\begin{commentblock}
\textbf{Q1. Justification of Practical Applicability:} The primary scope of this journal is automation in construction. How do the authors justify the practical deployment of a system that misclassifies roughly 50\% of the structural point cloud? The authors must provide a realistic discussion on what specific, limited engineering tasks (if any) an mIoU of 0.34 can actually support, rather than broadly claiming it ``reduces manual per-tunnel re-tuning.''
\end{commentblock}
\begin{responseblock}
We agree that broad deployment claims are not supported at mIoU $\approx$ 0.34. We have removed claims of ``reduces manual per-tunnel re-tuning'' throughout the revised manuscript and state explicitly that current accuracy does not support autonomous final inspection.

Within this scope, we discuss limited tasks that the framework \emph{may} support today, subject to human review:
\begin{itemize}[nosep]
  \item first-pass, bounded parameter adaptation before expert verification, producing a JSON parameter record and logged rationale;
  \item flagging when a target tunnel departs far from the calibrated reference anchor;
  \item shifting expert effort from repeated per-tunnel intervention toward upfront calibration and knowledge authoring (Discussion, Section~5.1).
\end{itemize}
We do \emph{not} claim readiness for structural health monitoring, deformation analysis, water-leakage detection, or Scan-to-BIM without human verification.
\end{responseblock}

\begin{commentblock}
\textbf{Q2. In-depth Error Analysis:} The manuscript currently lacks a qualitative and quantitative error analysis. Where are these massive errors concentrated? Are the LLMs failing to identify the correct boundaries? Is the 2D-to-3D reprojection introducing severe artifacts? Are background noises being classified as key blocks? The authors must provide visual error maps (e.g., False Positives vs.\ False Negatives) and analyze why the absolute metrics remain so stubbornly low despite LLM intervention.
\end{commentblock}
\begin{responseblock}
We agree. A new error-analysis subsection, Section~4.6 (Error analysis), has been added. It reports:
\begin{itemize}[nosep]
  \item FP/FN error maps on representative regular (tunnel 1-1) and complex (tunnel 4-1) examples (Fig.~11);
  \item a quantitative error-composition table over all ground-truth points (Table~9);
  \item quantified diagnostics for four structural limits of the fixed SAM4Tun labelling substrate (Table~10; Fig.~12): ring-level non-uniform point density, a moving K-anchor, fixed segment-offset template drift, and hard-coded walk direction.
\end{itemize}
The Discussion (Section~5.1) links these failure modes to the stated absolute performance ceiling and to the mechanism-level interpretation of the adaptation results. False positives (background classified as blocks) remain rare ($\leq$6\%); the dominant residual errors are class swaps from the fixed positional labelling rule, not 2D-to-3D reprojection artefacts or widespread noise misclassification as key blocks.
\end{responseblock}

\begin{commentblock}
\textbf{Q3. Comparison with State-of-the-Art (SOTA) Baselines:} The current comparison is strictly internal (LLM vs.\ Non-LLM vs.\ Static SAM4Tun). How does an absolute mIoU of 0.34 compare to existing, traditional geometric feature-engineering methods or standard supervised 3D deep learning models trained on similar datasets? If a standard, non-LLM baseline can achieve an mIoU of 0.65+, the justification for introducing the computational overhead and API costs of LLMs is severely weakened.
\end{commentblock}
\begin{responseblock}
We agree that internal comparisons alone are insufficient. We added Section~5.2 (Positioning relative to supervised and expert-tuned methods) and Table~11, which situates R4Tun against established paradigms on Seg2Tunnel:
\begin{itemize}[nosep]
  \item supervised 3D deep learning (PointNet++, RandLA-Net, voxel-based methods): mIoU $\approx$ 0.83--0.93, but requires per-point labels and GPU retraining;
  \item expert-tuned SAM4Tun: mIoU $>$0.90 on curated test rings, but requires per-case expert tuning;
  \item single-reference static SAM4Tun (this study): mIoU 0.15 without re-tuning.
\end{itemize}
R4Tun (mIoU 0.32--0.34) is well below these methods on absolute accuracy, and we do not claim otherwise. The comparison is on setup cost: R4Tun is the only entry in Table~11 that is simultaneously label-free, training-free, and expert-tuning-free while still adapting per tunnel. We therefore frame the paper as a mechanism study, not a ready-to-deploy replacement for supervised or expert-tuned segmentation.
\end{responseblock}

\begin{commentblock}
\textbf{Q4. The Ceiling of the Pipeline:} The authors briefly mention the ``single-reference design'' as a limitation. However, if the underlying SAM4Tun pipeline is inherently incapable of achieving high mIoU due to its 3D-to-2D projection and SAM-prompting mechanics, then using an LLM to tune its parameters is merely ``rearranging deck chairs on the Titanic.'' The authors need to clarify whether the low mIoU is a failure of the LLM's parameter selection or a fundamental limitation of the SAM4Tun architecture itself.
\end{commentblock}
\begin{responseblock}
We agree this distinction must be stated clearly. Section~4.6 identifies four structural constraints within the fixed SAM4Tun positional-labelling substrate:
\begin{enumerate}[nosep]
  \item \textbf{Non-uniform point density:} central rings are far denser than end rings (up to $19\times$ regular, $38\times$ complex); a single tunnel-level setup cannot span both.
  \item \textbf{Moving K-anchor:} staggered assembly shifts K along the arc; K-mislocated rings rise from 12\% (regular) to 64\% (complex), collapsing per-ring accuracy.
  \item \textbf{Fixed segment-offset template:} each block is placed at a fixed angular offset from K; the 6-segment template fits regular tunnels but fails on 7-segment complex geometry.
  \item \textbf{Hard-coded walk direction:} reversed-handedness rings produce mirror-image labels that parameter tuning cannot correct.
\end{enumerate}
The low absolute mIoU is therefore primarily a limitation of the current SAM4Tun open-source implementation. However, these constraints do not invalidate the underlying SAM4Tun architecture, nor the mechanism-level finding that structured context enables cross-LLM parameter adaptation. Instead, they highlight a clear future direction: improving the current SAM4Tun pipeline so that bounded LLM-based adaptation can be applied effectively.
\end{responseblock}

\begin{commentblock}
The integration of LLMs for context-aware parameter adaptation is interesting. The authors have done a good job structuring the experiments and ablation studies. Unfortunately, the actual numerical outcomes (mIoU $\sim$0.33, OA $\sim$0.53) reveal that the resulting system is riddled with segmentation errors and is currently unfit for practical, real-world tunnel inspection. Unless the authors can provide a compelling, evidence-based justification for why these low absolute metrics are acceptable, or demonstrate that the errors are benign for specific downstream tasks, the manuscript in its current form does not meet the practical impact threshold expected by the readership of \emph{Automation in Construction}.
\end{commentblock}
\begin{responseblock}
We agree that the system is not fit for autonomous real-world tunnel inspection at current accuracy. To address this concern, we have revised the manuscript to frame the contribution and to supply the evidence requested above. Summary of actions:
\begin{itemize}[nosep]
  \item \textbf{Reframed study scope:} repositioned as a mechanism-contribution study; R4Tun explores zero-label, no-retraining, no-expert-tuned self-adaptation and does not yet support autonomous final inspection.
  \item \textbf{Performance framing:} removed celebratory relative-gain language; discussion leads with absolute mIoU and OA and acknowledges the performance ceiling.
  \item \textbf{In-depth error analysis:} Section~4.6 (Figs.~11--12; Tables~9--10) decomposes errors and quantifies four structural bottlenecks of the fixed labelling substrate.
  \item \textbf{External benchmarking:} Section~5.2 and Table~11 situate R4Tun against supervised and expert-tuned tunnel-segmentation methods.
  \item \textbf{LLM vs.\ rule distinction:} Section~3.4 (agent design, context, CoT), Algorithm~1, and Section~3.5.2 (Non-LLM baseline) separate LLM reasoning from deterministic rule lookup.
  \item \textbf{Practical-utility claim constrained:} we explicitly limit claims to mechanism-level adaptation behavior and do not assert that the remaining segmentation errors are benign for downstream inspection tasks at this stage.
  \item \textbf{Limitations and future work:} Conclusion lists inherent limitations separately from future directions (dynamic labelling, multiple reference anchors, iterative self-correction).
\end{itemize}
\end{responseblock}

\section*{Reviewer 3}

\begin{commentblock}
The study objectives and general rationale are mostly clear. However, the role and generality of the LLM in the proposed framework require more precise framing. The current wording may suggest that R4Tun provides a general LLM-based tunnel segmentation capability, whereas the implemented workflow still relies on an expert-designed pipeline, including tunnel unfolding, denoising, image enhancement, SAM-based segmentation, and 3D reprojection. Please revise the motivation, especially around Fig.~1, to clarify that the main contribution is LLM-guided bounded parameter adaptation for an existing segmentation pipeline, rather than a generally capable LLM solution for tunnel segmentation.
\end{commentblock}
\begin{responseblock}
We agree. Fig.~1 and its caption have been revised to clarify that R4Tun extends a fixed expert-designed pipeline (SAM4Tun) via bounded parameter adaptation, rather than providing a general LLM-based tunnel segmentation capability. The Abstract, Introduction, and Conclusion use the same framing throughout.
\end{responseblock}

\begin{commentblock}
The method is generally described in sufficient detail, but some clarification would improve reader understanding. Additional explanation is needed to clarify how the LLM converts tunnel characteristics, reference parameters, intermediate pipeline states, and stage-specific knowledge documents into meaningful bounded parameter updates. Without this explanation, the specific role of LLM-based reasoning remains difficult to distinguish from deterministic rule-based parameter adjustment.
\end{commentblock}
\begin{responseblock}
We agree. Section~3.4 (Agent design, Context design, CoT design) and Algorithm~1 describe how each stage agent assembles memory, state, and knowledge into a prompt and runs a five-step CoT protocol. Section~3.5.2 and Appendix~E define the deterministic Non-LLM baseline (same knowledge encoded as explicit if--then lookups on characterisation fields).

Section~4.1, Table~4, and Fig.~9 compare Non-LLM vs.\ LLM: on regular tunnels, rules do not beat the static baseline; LLM ($m+s+k$) does; $m+s$ exceeds the rule baseline by 0.10--0.12 mIoU overall.

Appendix~C gives the denoising agent's memory, state, and knowledge inputs; Appendix~D gives a full five-step CoT trace and returned JSON for one tunnel.
\end{responseblock}

\begin{commentblock}
The manuscript would benefit from improving several existing figures and tables so that they are more self-contained.
\begin{itemize}[nosep]
  \item \textbf{Fig.~1:} The motivation in Fig.~1 should be revised to avoid overstating the role of the LLM. The figure should clarify that R4Tun reduces repeated deployment-stage expert retuning by reusing expert-defined reference parameters, parameter bounds, and knowledge documents, rather than removing expert knowledge from the workflow.
  \item \textbf{Fig.~2:} Some labels, such as B1 and K2, are difficult to understand from the figure alone. These labels should be defined directly in the figure or caption so that readers do not need to refer back to the main text.
  \item \textbf{Fig.~5:} The relationship between the three interpolation methods shown on the left and the workflow shown on the right is ambiguous. The figure should clarify whether these methods are used sequentially or presented as alternatives.
  \item \textbf{Table~7:} Some parameters, especially the underlined parameters, are not sufficiently explained. The authors should add short descriptions of what these parameters control and how they affect the corresponding stages of the segmentation pipeline.
\end{itemize}
\end{commentblock}
\begin{responseblock}
\textbf{Fig.~1:} revised to clarify that R4Tun reuses expert-defined reference parameters, parameter bounds, and knowledge documents, emphasising reduced deployment-stage retuning within a predefined pipeline rather than removal of expert knowledge.

\textbf{Fig.~2:} revised figure and caption define B1, K, and related labels directly.

\textbf{Fig.~5:} revised figure and caption clarify that the interpolation methods are sequential steps in one workflow.

\textbf{Table~7:} short descriptions added for the underlined parameters, explaining what each controls and how it affects the corresponding pipeline stage.
\end{responseblock}

\begin{commentblock}
The strengths of the study are present, but they could be emphasised with greater precision. The manuscript should more clearly state that the strength of R4Tun lies in integrating feature-engineered preprocessing, a foundation segmentation model, and LLM-guided bounded parameter adaptation. This framing would better highlight the actual contribution and avoid overstating the LLM as a general-purpose tunnel segmentation model.
\end{commentblock}
\begin{responseblock}
We have revised the Abstract, Introduction (contributions list), and Conclusion to state explicitly that R4Tun integrates feature-engineered preprocessing, a foundation segmentation model (SAM4Tun), and LLM-guided bounded parameter adaptation---not a standalone LLM segmentation capability.
\end{responseblock}

\begin{commentblock}
The limitations are generally stated, but the manuscript should more explicitly acknowledge that R4Tun does not eliminate the need for expert knowledge. Expert input is still required to define reference parameters, parameter bounds, and stage-specific knowledge documents.
\end{commentblock}
\begin{responseblock}
We have made this explicit in Section~5.3 (Limitations) and Section~6 (Conclusions). The revised manuscript presents the LLM as adapting within expert-defined constraints (reference parameters, bounds, and knowledge documents), not replacing them.
\end{responseblock}

\begin{commentblock}
The manuscript structure is generally clear, but the framing and figure explanations require revision to improve flow.
\end{commentblock}
\begin{responseblock}
We have revised the Abstract, Introduction, and Conclusion to reframe the scope of the framework. We have also revised the captions of Figs.~1, 2, and 5 to improve readability and manuscript flow.
\end{responseblock}

\bigskip
\noindent We believe the revised manuscript now presents R4Tun honestly as a controlled adaptation mechanism, reports absolute performance without overstated relative-gain framing, and provides the error decomposition and external positioning requested by Reviewer~2. We welcome any further guidance from the Editor and reviewers.

\end{document}
